\section{Discussion and Limitations}
\label{sec:discussion}

\subsection{A systems foundation, not a moral ranking}

The words \emph{facts}, \emph{trust}, and \emph{cooperation} carry moral and
political meanings. This paper does not rank participants by possession of
those virtues. It assigns operational responsibilities to a system: preserve
the reference, expose the warrant, and make the basis of coordinated action
inspectable. A participant may be honest while using a drifting contract; a
reputable institution may make an under-evidenced claim; competitors may
interoperate reliably without mutual affection.

\subsection{Facts are admitted, not total}

KFD-1 does not turn a database, repository, or ledger into reality. Every fact
cut has an observer, source, authority boundary, and omissions. Stability can
preserve a mistaken model. The recursive kernel matters because contrary
consequence can be retained without rewriting the earlier cut, making
correction possible. KFD-4 makes the perspective boundary more explicit; it
does not make all perspectives equally supported.

\subsection{Trust is bounded, not accumulated prestige}

Trust assessments are purpose-specific and revisable. A system may be trusted
to preserve bytes but not to interpret their value; an agent may be trusted to
run a bounded transformation but not to approve its deployment. Aggregating
these assessments into a universal trust score would discard the purpose,
evidence, and responsibility boundaries on which KFD-2 depends.

\subsection{Cooperation is not harmony or permissiveness}

Participants may disagree, compete, refuse, or operate under strict safety and
legal constraints. KFD-3 applies where a system asks reasoning participants to
coordinate. It requires participant-relevant value and constraints to be
inspectable; it does not guarantee equal bargaining power or eliminate
coercive conditions outside the system boundary. A complete evaluation must
therefore examine meaningful alternatives and affected non-participants, not
only interface transparency.

\subsection{Minimal roles are not a complete ontology}

The removal arguments establish that the three roles are non-redundant within
the proposed model. They do not prove that the English terms are uniquely
minimal or that every domain must expose them as first-class objects. A system
that preserves equivalent continuity through another decomposition is a valid
counterexample to vocabulary-level necessity. Security, incentives, law,
embodiment, resource constraints, institutional history, and domain semantics
remain outside the triad's complete account.

\subsection{Originating-ecosystem bias}

The model was developed inside the Kungfu ecosystem and is reflected in its
product, release, and agent-work designs. That coherence is implementation
evidence but also a source of overfitting. The cross-session case is a design
trace, not proof that the same decomposition minimizes cost in healthcare,
public administration, scientific collaboration, or other domains.
Independent adopters may find a narrower kernel, different boundaries, or
social practices that outperform formalization.

\subsection{Formalization has a cost}

Fact cuts, evidence bindings, responsibility assignments, and cooperation
records move reconstruction cost from participants into infrastructure. The
move is justified only where the recurring cost, consequence, or duration of
the work exceeds the cost of maintaining the mechanism. A short-lived local
task may need no formal KFD surface. The foundation is load-bearing where the
work must cross a boundary and still remain the same work.
